% --- CORRECCIÓN IMPORTANTE EN LA PRIMERA LÍNEA ---
% Agregamos "xcolor=table" para que funcione \rowcolor sin errores
\documentclass[aspectratio=169, 10pt, xcolor=table]{beamer}

% --- Configuración del Tema (estilo profesional como el ejemplo) ---
\usetheme{Madrid}
\usecolortheme{dolphin}

% --- Paquetes necesarios ---
\usepackage[utf8]{inputenc}
\usepackage[spanish]{babel}
\usepackage{amsmath, amssymb}
\usepackage{booktabs}
\usepackage{graphicx}
\usepackage{listings}
\usepackage{tikz}
\usetikzlibrary{arrows.meta, positioning, shapes.geometric}

% --- Ruta de Imágenes (crea la carpeta "imagenes" y coloca ahí tus archivos) ---
\graphicspath{{./imagenes/}}

% --- Estilo para código Python ---
\definecolor{codegreen}{rgb}{0,0.6,0}
\definecolor{codegray}{rgb}{0.5,0.5,0.5}
\definecolor{codepurple}{rgb}{0.58,0,0.82}
\definecolor{backcolour}{rgb}{0.95,0.95,0.92}
\lstset{
    language=Python,
    backgroundcolor=\color{backcolour},
    commentstyle=\color{codegreen},
    keywordstyle=\color{blue},
    numberstyle=\tiny\color{codegray},
    stringstyle=\color{codepurple},
    basicstyle=\ttfamily\scriptsize,
    breaklines=true,
    frame=single,
    showstringspaces=false
}

% --- Pie de página personalizado ---
\makeatletter
\setbeamertemplate{footline}{
  \leavevmode%
  \hbox{%
  \begin{beamercolorbox}[wd=.333333\paperwidth,ht=2.25ex,dp=1ex,center]{author in head/foot}%
    \usebeamerfont{author in head/foot}\insertshortauthor
  \end{beamercolorbox}%
  \begin{beamercolorbox}[wd=.333333\paperwidth,ht=2.25ex,dp=1ex,center]{title in head/foot}%
    \usebeamerfont{title in head/foot}Sesión 2
  \end{beamercolorbox}%
  \begin{beamercolorbox}[wd=.333333\paperwidth,ht=2.25ex,dp=1ex,right]{date in head/foot}%
    \usebeamerfont{date in head/foot}Machine Learning \hfill \insertframenumber/\inserttotalframenumber \hspace*{0.3cm}
  \end{beamercolorbox}}%
  \vskip0pt%
}
\makeatother

% --- Datos de la portada ---
\title[SESIÓN 2]{\textbf{Sesión 2}}  
\subtitle{Análisis Exploratorio y Feature Engineering}
\author[Prof. C. Sánchez]{Carlos César Sánchez Coronel}
\institute{\textbf{MACHINE LEARNING}} 
\date{}

% --- Eliminar la barra de navegación (opcional) ---
\setbeamertemplate{navigation symbols}{}

% --- Índice al inicio de cada sección ---
\AtBeginSection[]{
  \begin{frame}
    \frametitle{Contenido}
    \tableofcontents[currentsection]
  \end{frame}
}

% ============================================================================
% INICIO DEL DOCUMENTO
% ============================================================================
\begin{document}

% --- Portada ---
\begin{frame}
    \titlepage
\end{frame}

% --- Índice general ---
\begin{frame}{Contenido}
    \tableofcontents
\end{frame}

% ============================================================================
% SECCIÓN 1: INTRODUCCIÓN Y ESTADÍSTICAS ROBUSTAS
% ============================================================================
\section{Introducción: El arte detectivesco del EDA}

\begin{frame}{EDA como proceso detectivesco}
    \begin{block}{Filosofía de John Tukey}
        “Escuchar a los datos” mediante visualización y estadísticas robustas antes de modelar.
    \end{block}
    \begin{exampleblock}{Objetivo del Feature Engineering}
        Traducir hallazgos en características matemáticas que maximicen el rendimiento de los algoritmos.
    \end{exampleblock}
    \begin{itemize}
        \item Juntos, EDA y Feature Engineering determinan el límite superior del rendimiento de cualquier modelo.
    \end{itemize}
    % IMAGEN: Ilustración de detective con lupa sobre datos
\end{frame}

\section{Estadísticas robustas vs clásicas}

\begin{frame}{Estadísticas descriptivas robustas}
    % Basado en Pregunta 1 del solucionario
    \begin{columns}[t]
        \column{0.5\textwidth}
        \begin{alertblock}{Clásicas (sensibles a outliers)}
            \begin{itemize}
                \item Media: $\bar{x} = \frac{1}{n}\sum x_i$
                \item Desviación estándar: $s = \sqrt{\frac{1}{n-1}\sum (x_i-\bar{x})^2}$
            \end{itemize}
        \end{alertblock}
        \column{0.5\textwidth}
        \begin{exampleblock}{Robustas (resistentes)}
            \begin{itemize}
                \item Mediana: $Median = x_{(n/2)}$
                \item MAD: $MAD = median(|x_i - median(x)|)$
            \end{itemize}
        \end{exampleblock}
    \end{columns}
    \vspace{0.3cm}
    \begin{block}{¿Cuándo usar robustas?}
        Cuando hay asimetría o outliers. Ej: salarios (media > mediana $\rightarrow$ sesgo derecho).
    \end{block}
    % Basado en Pregunta 1
\end{frame}

\begin{frame}{Detección de outliers con z-score modificado}
    \begin{block}{Z-score modificado (basado en MAD)}
        \[
        M_i = \frac{0.6745(x_i - \text{mediana}(x))}{\text{MAD}}
        \]
        \begin{itemize}
            \item Factor 0.6745 hace consistente la MAD con la desviación estándar en datos normales.
            \item Umbral típico: $|M_i| > 3.5$ indica outlier.
        \end{itemize}
    \end{block}
    % Basado en Pregunta 1c
    % IMAGEN: Diagrama de caja con outliers
\end{frame}

% ============================================================================
% SECCIÓN 2: VISUALIZACIONES AVANZADAS
% ============================================================================
\section{Visualizaciones avanzadas}

\begin{frame}{Histograma y KDE}
    % Basado en Pregunta 2
    \begin{columns}[t]
        \column{0.5\textwidth}
        \begin{block}{Histograma}
            \begin{itemize}
                \item Ancho de bin óptimo (Freedman-Diaconis): 
                \[
                \text{ancho} = 2 \cdot \text{IQR} \cdot n^{-1/3}
                \]
                \item Útil para distribuciones con colas pesadas.
            \end{itemize}
        \end{block}
        \column{0.5\textwidth}
        \begin{block}{KDE (Kernel Density Estimation)}
            \[
            \hat{f}(x) = \frac{1}{nh}\sum_{i=1}^n K\left(\frac{x-x_i}{h}\right)
            \]
            \begin{itemize}
                \item $h$: ancho de banda (controla suavizado).
                \item $K$: kernel (gaussiano por defecto).
            \end{itemize}
        \end{block}
    \end{columns}
    % IMAGEN: Comparación histograma vs KDE
\end{frame}

\begin{frame}{ECDF y Mapas de calor de correlación}
    \begin{columns}[t]
        \column{0.5\textwidth}
        \begin{block}{ECDF (Empirical CDF)}
            \[
            F_n(x) = \frac{1}{n}\sum_{i=1}^n \mathbf{1}_{x_i \le x}
            \]
            Muestra la proporción acumulada. Útil para comparar distribuciones.
        \end{block}
        \column{0.5\textwidth}
        \begin{block}{Correlación de Pearson y Spearman}
            Pearson: relación lineal. \\
            Spearman: relación monótona (basada en rangos). \\
            \vspace{0.2cm}
            \textbf{Mapa de calor:} visualiza multicolinealidad.
        \end{block}
    \end{columns}
    % Basado en Pregunta 3
    % IMAGEN: Mapa de calor de correlaciones
\end{frame}

% ============================================================================
% SECCIÓN 3: DETECCIÓN DE OUTLIERS AVANZADA
% ============================================================================
\section{Detección de outliers avanzada}

\begin{frame}{Métodos univariados: IQR y z-score modificado}
    % Basado en Pregunta 4
    \begin{block}{Regla del IQR}
        \begin{itemize}
            \item IQR = Q3 – Q1
            \item Límites: $[Q1 - 1.5\cdot IQR,\ Q3 + 1.5\cdot IQR]$
            \item Limitación: asume simetría; no detecta outliers contextuales.
        \end{itemize}
    \end{block}
    \begin{exampleblock}{Ejemplo práctico: detección de fraudes}
        En montos de transacciones, aplicar IQR por separado puede no detectar un monto normal en horario inusual. Para eso necesitamos métodos multivariados.
    \end{exampleblock}
\end{frame}

\begin{frame}{Métodos multivariados (I)}
    % Basado en Pregunta 5
    \begin{block}{Distancia de Mahalanobis}
        \[
        D_M(x) = \sqrt{(x-\mu)^T \Sigma^{-1}(x-\mu)}
        \]
        \begin{itemize}
            \item Considera correlaciones entre variables.
            \item Asume normalidad multivariante.
            \item Se considera outlier si $D_M^2 > \chi^2_{p,0.975}$.
        \end{itemize}
    \end{block}
    % IMAGEN: Elipse de Mahalanobis vs euclidiana
\end{frame}

\begin{frame}{Métodos multivariados (II)}
    % Basado en Preguntas 5 y 6
    \begin{columns}[t]
        \column{0.5\textwidth}
        \begin{block}{LOF (Local Outlier Factor)}
            Compara densidad local con vecinos.
            \[
            LOF(p) = \frac{\frac{1}{k}\sum_{o\in N_k(p)} lrd(o)}{lrd(p)}
            \]
            LOF $\gg 1$ indica outlier local.
        \end{block}
        \column{0.5\textwidth}
        \begin{block}{Isolation Forest}
            Aísla puntos con particiones aleatorias.
            \begin{itemize}
                \item Outliers requieren menos particiones.
                \item Puntuación de anomalía: $s(x)=2^{-\frac{E(h(x))}{c(n)}}$.
                \item $s \approx 1$ $\rightarrow$ anomalía.
            \end{itemize}
        \end{block}
    \end{columns}
    % IMAGEN: Ejemplo de LOF e Isolation Forest
\end{frame}

% ============================================================================
% SECCIÓN 4: MANEJO DE DATOS FALTANTES
% ============================================================================
\section{Manejo de datos faltantes}

\begin{frame}{Mecanismos de ausencia}
    % Basado en Pregunta 7
    \begin{columns}[t]
        \column{0.33\textwidth}
        \begin{block}{MCAR}
            Falta completamente al azar (ej. error de transcripción).
        \end{block}
        \column{0.33\textwidth}
        \begin{block}{MAR}
            Falta al azar condicionado a otras variables (ej. mujeres no responden ingresos).
        \end{block}
        \column{0.33\textwidth}
        \begin{block}{MNAR}
            Falta no al azar (depende del valor faltante mismo, ej. altos ingresos no declaran).
        \end{block}
    \end{columns}
    \vspace{0.3cm}
    \begin{alertblock}{Importancia}
        Identificar el mecanismo ayuda a elegir la imputación adecuada. Ignorar MNAR sesga estimaciones.
    \end{alertblock}
\end{frame}

\begin{frame}{Imputación simple vs múltiple}
    % Basado en Pregunta 8
    \begin{columns}[t]
        \column{0.5\textwidth}
        \begin{block}{Imputación simple}
            \begin{itemize}
                \item Media/mediana: rápido, subestima varianza.
                \item Regresión: usa otras variables.
                \item kNN: promedio de vecinos.
            \end{itemize}
        \end{block}
        \column{0.5\textwidth}
        \begin{block}{Imputación múltiple (MICE)}
            \begin{itemize}
                \item Genera $m$ conjuntos (ej. 5-10) mediante modelos iterativos.
                \item Combina resultados con reglas de Rubin.
                \item Incorpora incertidumbre de la imputación.
            \end{itemize}
        \end{block}
    \end{columns}
    % IMAGEN: Esquema de MICE
\end{frame}

% ============================================================================
% SECCIÓN 5: FEATURE ENGINEERING AVANZADO
% ============================================================================
\section{Feature Engineering Avanzado}

\subsection{Transformaciones de variables}

\begin{frame}{Normalización, estandarización y escalado robusto}
    % Basado en Pregunta 9
    \begin{columns}[t]
        \column{0.33\textwidth}
        \begin{block}{Normalización (Min-Max)}
            $x' = \frac{x - x_{min}}{x_{max} - x_{min}} \in [0,1]$
        \end{block}
        \column{0.33\textwidth}
        \begin{block}{Estandarización (Z-score)}
            $x' = \frac{x - \mu}{\sigma}$
        \end{block}
        \column{0.33\textwidth}
        \begin{block}{Escalado robusto}
            $x' = \frac{x - \text{mediana}}{IQR}$
        \end{block}
    \end{columns}
    \vspace{0.2cm}
    \begin{itemize}
        \item SVM y PCA requieren estandarización.
        \item Escalado robusto: preferible cuando hay outliers.
    \end{itemize}
\end{frame}

\begin{frame}{Transformaciones para asimetría}
    % Basado en Pregunta 10
    \begin{columns}[t]
        \column{0.5\textwidth}
        \begin{block}{Logarítmica}
            $x' = \log(x)$ para datos positivos con cola derecha.
        \end{block}
        \begin{block}{Box-Cox}
            \[
            x' = \begin{cases}
                \frac{x^\lambda - 1}{\lambda} & \lambda \neq 0 \\
                \log(x) & \lambda = 0
            \end{cases}
            \]
            Requiere $x>0$.
        \end{block}
        \column{0.5\textwidth}
        \begin{block}{Yeo-Johnson}
            Extensión de Box-Cox que acepta valores negativos.
        \end{block}
    \end{columns}
    \vspace{0.2cm}
    % Basado en Pregunta 10
\end{frame}

\subsection{Creación de nuevas características}

\begin{frame}{Términos polinomiales e interacciones}
    % Basado en Pregunta 11
    \begin{block}{Polinomiales}
        $x^2, x^3$ capturan relaciones no lineales. Cuidado con overfitting.
    \end{block}
    \begin{block}{Interacciones}
        $x_i \cdot x_j$ modelan sinergias. Ejemplo: en precios de viviendas, la interacción entre metros cuadrados y número de habitaciones puede ser más informativa.
    \end{block}
    % IMAGEN: Superficie de respuesta con interacción
\end{frame}

\begin{frame}{Binning y discretización}
    % Basado en Pregunta 12
    \begin{columns}[t]
        \column{0.5\textwidth}
        \begin{block}{Igual ancho}
            Divide el rango en intervalos fijos. Sensible a outliers.
        \end{block}
        \column{0.5\textwidth}
        \begin{block}{Igual frecuencia}
            Intervalos con igual número de observaciones. Mejor para distribuciones asimétricas.
        \end{block}
    \end{columns}
    \vspace{0.2cm}
    \begin{itemize}
        \item Discretizar pierde información, pero puede mejorar la robustez.
    \end{itemize}
\end{frame}

\begin{frame}{Variables temporales y agregadas}
    % Basado en Preguntas 13 y 14
    \begin{columns}[t]
        \column{0.5\textwidth}
        \begin{block}{Rezagos y ventanas móviles}
            \begin{itemize}
                \item $y_{t-1}, y_{t-2}$: dependencia temporal.
                \item Media móvil: $MA_t = \frac{1}{k}\sum_{i=0}^{k-1} y_{t-i}$
                \item Desviación móvil: volatilidad.
            \end{itemize}
        \end{block}
        \column{0.5\textwidth}
        \begin{block}{Codicación cíclica (hora, día)}
            \[
            \sin\left(\frac{2\pi \cdot \text{hora}}{24}\right),\quad \cos\left(\frac{2\pi \cdot \text{hora}}{24}\right)
            \]
            Evita la discontinuidad entre 23h y 0h.
        \end{block}
    \end{columns}
\end{frame}

\begin{frame}{Variables agregadas por grupo}
    % Basado en Pregunta 14
    \begin{block}{Agregaciones por cliente/usuario}
        \begin{itemize}
            \item Media, desviación, conteo, min, max, etc.
            \item Ejemplo: en detección de fraude, frecuencia de transacciones en últimas 24h, monto promedio, etc.
        \end{itemize}
    \end{block}
    \begin{alertblock}{¡Cuidado con el data leakage!}
        Usar ventanas temporales: solo datos anteriores al momento de predicción.
    \end{alertblock}
\end{frame}

\subsection{Codificación de variables categóricas}

\begin{frame}{Codificación de alta cardinalidad}
    % Basado en Pregunta 15
    \begin{columns}[t]
        \column{0.33\textwidth}
        \begin{block}{One-hot encoding}
            Crea muchas columnas. No escalable a cientos de categorías.
        \end{block}
        \column{0.33\textwidth}
        \begin{block}{Target encoding}
            Reemplaza categoría por media de la variable objetivo. ¡Requiere validación cruzada para evitar leakage!
        \end{block}
        \column{0.33\textwidth}
        \begin{block}{Frequency encoding}
            Reemplaza por frecuencia de aparición. Útil en NLP.
        \end{block}
    \end{columns}
    \vspace{0.2cm}
    \begin{block}{Weight of Evidence (WoE) para clasificación binaria}
        $WoE = \ln\left(\frac{\%eventos}{\%no eventos}\right)$
    \end{block}
\end{frame}

\subsection{Selección de características}

\begin{frame}{Métodos de selección}
    % Basado en Preguntas 16 y 17
    \begin{columns}[t]
        \column{0.33\textwidth}
        \begin{block}{Filtro}
            \begin{itemize}
                \item Correlación
                \item ANOVA (F-test)
                \item Información mutua: $I(X;Y)=\sum p(x,y)\log\frac{p(x,y)}{p(x)p(y)}$
            \end{itemize}
        \end{block}
        \column{0.33\textwidth}
        \begin{block}{Wrapper}
            \begin{itemize}
                \item Selección hacia adelante
                \item Eliminación hacia atrás
                \item Costosos pero pueden encontrar interacciones.
            \end{itemize}
        \end{block}
        \column{0.33\textwidth}
        \begin{block}{Embedded}
            \begin{itemize}
                \item Importancia de árboles (Random Forest)
                \item Regularización L1 (Lasso): $\min \sum (y-X\beta)^2 + \lambda\sum|\beta_j|$
            \end{itemize}
        \end{block}
    \end{columns}
\end{frame}

% ============================================================================
% SECCIÓN 6: PIPELINES Y AUTOMATIZACIÓN
% ============================================================================
\section{Pipelines y automatización}

\begin{frame}{Pipelines en scikit-learn}
    % Basado en Pregunta 18
    \begin{block}{Ventajas}
        \begin{itemize}
            \item Evitan data leakage (transformaciones solo en entrenamiento).
            \item Facilitan validación cruzada y búsqueda de hiperparámetros.
            \item Garantizan reproducibilidad.
        \end{itemize}
    \end{block}
    \begin{exampleblock}{Ejemplo}
        Pipeline con pasos: imputación, escalado, selección de características y modelo final.
    \end{exampleblock}
    % IMAGEN: Esquema de pipeline
\end{frame}

\begin{frame}{Feature engineering automático}
    % Basado en Pregunta 19
    \begin{block}{Deep Feature Synthesis (Featuretools)}
        Genera automáticamente características a partir de relaciones entre tablas (agregaciones, transformaciones).
    \end{block}
    \begin{itemize}
        \item Útil para datos relacionales (clientes, transacciones, productos).
        \item Precaución: limitar profundidad y usar selección posterior para evitar cientos de características irrelevantes.
    \end{itemize}
\end{frame}

% ============================================================================
% SECCIÓN 7: CASOS DE USO INTEGRADOS
% ============================================================================
\section{Casos de uso integrados}

\begin{frame}{Caso 1: Predicción de precios de viviendas}
    % Basado en Pregunta 19
    \begin{itemize}
        \item \textbf{EDA}: Detectar asimetría en precio (aplicar log), outliers en área (IQR), correlaciones altas (eliminar redundantes).
        \item \textbf{Feature engineering}: Interacciones (habitaciones × baños), variables temporales (año de construcción transformado con Box-Cox), target encoding de ubicación, escalado robusto.
    \end{itemize}
    % IMAGEN: Boxplots por ubicación
\end{frame}

\begin{frame}{Caso 2: Detección de fraude financiero}
    % Basado en Pregunta 20 (el gran caso integrador)
    \begin{columns}[t]
        \column{0.5\textwidth}
        \textbf{EDA:}
        \begin{itemize}
            \item Histogramas de montos (log scale)
            \item Boxplots por hora y día
            \item Mapas de calor de transacciones por ubicación
        \end{itemize}
        \textbf{Feature engineering específico:}
        \begin{itemize}
            \item Log del monto
            \item Variables temporales cíclicas
            \item Agregadas por usuario (frecuencia, monto promedio, desviación en últimas 24h)
            \item Distancia geográfica y temporal entre transacciones consecutivas
            \item Puntuación de anomalía de Isolation Forest como característica
        \end{itemize}
        \column{0.5\textwidth}
        \textbf{Codificación:} Target encoding con validación cruzada para tipo de comercio y ubicación. \\
        \textbf{Manejo de desbalance:} SMOTE o pesos en pérdida. \\
        \textbf{Pipeline:} Usar ventanas temporales para evitar leakage. \\
        \textbf{Métricas:} AUC-ROC, precisión, recall, y en producción A/B test con tasa de fraude detectado.
    \end{columns}
\end{frame}

% ============================================================================
% SECCIÓN 8: CONCLUSIÓN
% ============================================================================
\section{Conclusión}

\begin{frame}{Resumen de la sesión}
    \begin{exampleblock}{Lo que hemos visto}
        \begin{itemize}
            \item Estadísticas robustas vs clásicas (media/mediana, MAD).
            \item Visualizaciones avanzadas: KDE, ECDF, mapas de calor.
            \item Detección de outliers: IQR, Mahalanobis, LOF, Isolation Forest.
            \item Manejo de datos faltantes: MCAR/MAR/MNAR, MICE.
            \item Transformaciones: estandarización, Box-Cox, Yeo-Johnson.
            \item Creación de características: polinomiales, interacciones, binning, variables temporales y agregadas.
            \item Codificación de categóricas: target encoding, WoE.
            \item Selección de características: filtro, wrapper, embedded.
            \item Pipelines y automatización (Featuretools).
            \item Casos prácticos: viviendas y fraude.
        \end{itemize}
    \end{exampleblock}
    \vspace{0.2cm}

\end{frame}

% --- Diapositiva final ---
\begin{frame}
    \centering
    \Huge \textbf{¡Gracias!}

\end{frame}

\end{document}